% Generado por scripts/tablas.py a partir de Practica3_pendulos_acoplados.xlsx. No editar a mano.

\begin{table}[ht]
\centering
\caption{Periodos del modo en fase y frecuencia propia de cada medida. La última fila recoge el periodo medio y la frecuencia $\omega_0=2\pi/\mean{T}_0$ que empleamos en adelante.}
\label{tab:omega0}
\renewcommand{\arraystretch}{1.22}
\begin{tabular}{@{}l c c c@{}}
\toprule
Longitud & Péndulo & $T_0$ (\si{\second}) & $2\pi/T_0$ (\si{\radian\per\second}) \\
\midrule
$\ell_1=\SI{25}{\centi\meter}$ & 1 & \num{2.0056} & \num{3.1328} \\
$\ell_1=\SI{25}{\centi\meter}$ & 2 & \num{1.9904} & \num{3.1567} \\
$\ell_2=\SI{35}{\centi\meter}$ & 1 & \num{1.9912} & \num{3.1555} \\
$\ell_2=\SI{35}{\centi\meter}$ & 2 & \num{1.9968} & \num{3.1466} \\
\midrule
\multicolumn{2}{@{}l}{\scshape Media} & \num{1.9960} & \num{3.1479} \\
\bottomrule
\end{tabular}
\end{table}
